% context: markscheme (official solutions) for 2-5CW_Falling-objects
% topic: free fall with a = g = 9.80 m/s^2; dropped from rest, distance grows as t^2,
%        a ball thrown upward, v = 0 but a = -g at the top, air resistance
% convention: answers in bold; "=" joins exact equalities, "\approx" introduces a
%        rounded value; each solution lists the knowns, then names the equation used
\documentclass[12pt, twoside]{article}
%\documentclass[12pt, twoside]{article}
\usepackage[letterpaper, margin=1in, headsep=0.2in]{geometry}
\setlength{\headheight}{0.6in}
%\usepackage[english]{babel}
\usepackage[utf8]{inputenc}
\usepackage{microtype}
\usepackage{amsmath}
\usepackage{amssymb}
%\usepackage{amsfonts}
\usepackage[nomessages]{fp} %\FPeval{\var-name}{2*sin(pi/6)}
\usepackage{siunitx} %units in math. eg 20\milli\meter
\usepackage{yhmath} % for arcs, overparenth command
\usepackage{tikz} %graphics
\usetikzlibrary{quotes, angles, arrows, arrows.meta}
\usepackage{graphicx} %consider setting \graphicspath{{images/}}
\usepackage{parskip} %no paragraph indent
\usepackage{enumitem}
\usepackage{multicol}
\usepackage{venndiagram}

\usepackage{fancyhdr}
\pagestyle{fancy}
\fancyhf{}
\renewcommand{\headrulewidth}{0pt} % disable the underline of the header
\raggedbottom
\hfuzz=2mm %suppresses overfull box warnings

\usepackage{hyperref}
\usepackage{float}
\setlength{\parskip}{0.2\baselineskip}

\title{Physics}
\author{Chris Huson}
\date{October 2026}

\fancyhead[LE]{\thepage}
\fancyhead[RO]{\thepage}
\fancyhead[LO]{La Scuola d'Italia / Huson / Physics \\* 30 October 2026}

\begin{document}

\subsubsection*{2.5 Classwork: Falling Objects Solutions}

\begin{enumerate}[itemsep=0.5cm]

\item Dropped ball. Knowns: $v_0 = 0$, $y_0 = 0$, $a = +9.80$ m/s$^2$ (down
  positive).
  \begin{enumerate}[itemsep=0.3cm]
    \item \renewcommand{\arraystretch}{1.4}
      \begin{tabular}[t]{| c | c | c | c | c |}
        \hline
        $t$ (s) & 0 & 1.00 & 2.00 & 3.00 \\ \hline
        $v$ (m/s) & \textbf{0} & \textbf{9.80} & \textbf{19.6} & \textbf{29.4}
          \\ \hline
        $y$ (m) & \textbf{0} & \textbf{4.90} & \textbf{19.6} & \textbf{44.1}
          \\ \hline
      \end{tabular}

      \emph{$v = at$ adds 9.80 m/s each second. $y = \frac{1}{2}at^2$: for example
      $\frac{1}{2}(9.80~\text{m/s}^2)(2.00~\text{s})^2 = 4.90 \times 4.00 = 19.6$ m.
      This is Giancoli's Example 2--10.}
    \item $\dfrac{19.6}{4.90} = \mathbf{4}$ \textbf{times as far}, and
      $\dfrac{44.1}{4.90} = \mathbf{9}$ \textbf{times as far}.

      \textbf{The distance grows as the square of the time}: twice the time gives
      $2^2 = 4$ times the distance, three times the time gives $3^2 = 9$ times.

      \emph{This is Galileo's result, $d \propto t^2$. Also accept ``the ball falls
      farther each second'': 4.90 m, then 14.7 m, then 24.5 m, in the ratio
      1 : 3 : 5.}
  \end{enumerate}

\item Leaning Tower of Pisa. Knowns: $v_0 = 0$, $y - y_0 = 55.86$ m,
  $a = +9.80$ m/s$^2$ (down positive).
  \begin{enumerate}[itemsep=0.3cm]
    \item Want $t$. Use $y = y_0 + v_0 t + \frac{1}{2}at^2$ with $v_0 = 0$:

      $55.86~\text{m} = \tfrac{1}{2}(9.80~\text{m/s}^2)\,t^2$, \quad so
      $t^2 = \dfrac{2(55.86~\text{m})}{9.80~\text{m/s}^2} = 11.4~\text{s}^2$

      $t = \sqrt{11.4~\text{s}^2} = 3.3763\ldots \approx \mathbf{3.38~\textbf{s}}$

      \emph{A common slip is to forget the $\frac{1}{2}$, giving $t^2 = 5.7$ and
      $t \approx 2.39$ s, or to stop at $t^2 = 11.4$.}
    \item Want $v$. Use $v = v_0 + at$, carrying the unrounded time:

      $v = 0 + (9.80~\text{m/s}^2)(3.3763\ldots~\text{s}) = 33.088\ldots
      \approx \mathbf{33.1~\textbf{m/s}}$ downward

      \emph{Also accept $v^2 = v_0^2 + 2a(y - y_0) = 2(9.80)(55.86) = 1094.856$, so
      $v \approx 33.1$ m/s; this route does not need part (a). About 119 km/h.}
    \item \textbf{They land at the same time.} Without air resistance, every object
      falls with the same acceleration $g$, whatever its mass.

      \emph{Before Galileo, it was believed that a heavier object falls faster. For
      two dense balls the air makes almost no difference; it matters for a sheet
      of paper or a feather (Giancoli Figs.\ 2--20 and 2--21).}
  \end{enumerate}

\newpage

\item Ball thrown upward. Knowns: $v_0 = +14.7$ m/s, $a = -9.80$ m/s$^2$, $y_0 = 0$
  at the hand (up positive).
  \begin{enumerate}[itemsep=0.3cm]
    \item \textbf{Velocity $\mathbf{v = 0}$} for an instant. \textbf{Acceleration
      $\mathbf{a = -9.80}$ m/s$^\mathbf{2}$}, the same as everywhere else in the
      flight.

      \emph{The common wrong answer is $a = 0$ at the top. Gravity does not switch
      off: if the acceleration were zero there, the velocity would stay zero and the
      ball would hang in the air (Giancoli Example 2--14).}
    \item Want $t$ when $v = 0$. Use $v = v_0 + at$, so $t = \dfrac{v - v_0}{a}$:

      $t = \dfrac{0 - 14.7~\text{m/s}}{-9.80~\text{m/s}^2} = \mathbf{1.50~\textbf{s}}$
    \item Want $y - y_0$. Use $\bar{v} = \dfrac{v + v_0}{2} = \dfrac{0 + 14.7}{2}
      = 7.35$ m/s:

      $y - y_0 = \bar{v}\,t = 7.35~\text{m/s} \times 1.50~\text{s} = 11.025
      \approx \mathbf{11.0~\textbf{m}}$

      \emph{Also accept $v^2 = v_0^2 + 2a(y - y_0)$:
      $y - y_0 = \frac{0 - (14.7)^2}{2(-9.80)} = \frac{-216.09}{-19.6} = 11.025$ m.}
    \item \textbf{3.00 s.} The fall back down takes as long as the rise:
      $2 \times 1.50~\text{s} = 3.00$ s.

      $v = v_0 + at = 14.7~\text{m/s} + (-9.80~\text{m/s}^2)(3.00~\text{s})
      = \mathbf{-14.7~\textbf{m/s}}$, that is, 14.7 m/s \textbf{downward}.

      \emph{The ball returns at the speed it was thrown, in the opposite direction;
      the minus sign carries the direction. Also accept setting $y = 0$ in
      $y = v_0 t + \frac{1}{2}at^2$: $0 = (14.7 - 4.90\,t)\,t$, so $t = 0$ or
      $t = 3.00$ s (Giancoli Example 2--13).}
  \end{enumerate}

\item Raindrop. Knowns: $v_0 = 0$, $y - y_0 = 1000$ m, $a = +9.80$ m/s$^2$ (down
  positive).
  \begin{enumerate}[itemsep=0.3cm]
    \item Want $v$; no time is given. Use $v^2 = v_0^2 + 2a(y - y_0)$:

      $v^2 = 0 + 2(9.80~\text{m/s}^2)(1000~\text{m}) = 19{,}600~\text{m}^2/\text{s}^2$,
      \quad so $v = \mathbf{140~\textbf{m/s}}$

      \emph{About 500 km/h. Also accept finding $t = \sqrt{2(1000)/9.80}
      \approx 14.3$ s first, then $v = at \approx 140$ m/s.}
    \item \textbf{Air resistance.} The air pushes back on the drop, and the push
      grows as the drop goes faster, until it balances gravity. After that the drop
      stops speeding up and falls at a steady \emph{terminal velocity}, about 9 m/s.

      \emph{Air resistance matters most for small, light objects with a large
      surface for their weight --- like the flat sheet of paper in Giancoli
      Fig.\ 2--20. A raindrop reaches terminal velocity after falling only a few
      meters, so ignoring air here gives an answer more than 15 times too large.}
  \end{enumerate}

\end{enumerate}
\end{document}
